\textbf{Phần 1. Trắc nghiệm lựa chọn}
\setcounterpageref{ex}{0}

\begin{ex} %[4D3V11]
Lớp đơn vị gồm các hàng nào?
\choice
{Hàng chục, hàng trăm, hàng nghìn}
{\True Hàng đơn vị, hàng chục, hàng trăm}
{Hàng nghìn, hàng chục nghìn, hàng trăm nghìn}
{Hàng triệu, hàng chục triệu, hàng trăm triệu}
\loigiai{
Lớp đơn vị gồm: hàng đơn vị, hàng chục, hàng trăm.\\
\textbf{Đáp án của câu hỏi là B}
}
\end{ex}

\begin{ex} %[4D3V11]
Lớp nghìn gồm các hàng nào?
\choice
{Hàng đơn vị, hàng chục, hàng trăm}
{Hàng chục nghìn, hàng trăm nghìn, hàng triệu}
{Hàng nghìn, hàng triệu, hàng tỉ}
{\True Hàng nghìn, hàng chục nghìn, hàng trăm nghìn}
\loigiai{
Lớp nghìn gồm: hàng nghìn, hàng chục nghìn, hàng trăm nghìn.\\
\textbf{Đáp án của câu hỏi là D}
}
\end{ex}

\begin{ex} %[4D3V11]
Trong số $543\,210$, các chữ số thuộc lớp nghìn là
\choice
{\True $5; 4; 3$}
{$2; 1; 0$}
{$4; 3; 2$}
{$5; 4; 2$}
\loigiai{
Lớp nghìn gồm 3 chữ số: $5, 4, 3$.\\
\textbf{Đáp án của câu hỏi là A}
}
\end{ex}

\begin{ex} %[4D3V11]
Trong số $789\,123$, các chữ số thuộc lớp đơn vị là
\choice
{$7; 8; 9$}
{$9; 1; 2$}
{\True $1; 2; 3$}
{$8; 9; 1$}
\loigiai{
Lớp đơn vị gồm các chữ số: $1, 2, 3$.\\
\textbf{Đáp án của câu hỏi là C}
}
\end{ex}

\begin{ex} %[4D3V11]
Số nào sau đây có chữ số hàng chục nghìn bằng 8?
\choice
{$874\,215$}
{\True $382\,540$}
{$148\,900$}
{$520\,840$}
\loigiai{
Số $382\,540$ có chữ số 8 ở hàng chục nghìn.\\
\textbf{Đáp án của câu hỏi là B}
}
\end{ex}

\begin{ex} %[4D3V11]
Mỗi lớp trong hệ thống số thập phân gồm có bao nhiêu hàng?
\choice
{$2$ hàng}
{$4$ hàng}
{$5$ hàng}
{\True $3$ hàng}
\loigiai{
Mỗi lớp gồm 3 hàng liên tiếp.\\
\textbf{Đáp án của câu hỏi là D}
}
\end{ex}

\begin{ex} %[4D3V11]
Số nào sau đây có lớp nghìn gồm các chữ số $6; 0; 5$ và lớp đơn vị gồm các chữ số $0; 4; 8$?
\choice
{\True $605\,048$}
{$650\,480$}
{$605\,480$}
{$650\,048$}
\loigiai{
Số đó là $605\,048$.\\
\textbf{Đáp án của câu hỏi là A}
}
\end{ex}

\begin{ex} %[4D3V11]
Trong số $450\,325$, chữ số $5$ ở hàng chục nghìn có giá trị gấp bao nhiêu lần chữ số $5$ ở hàng đơn vị?
\choice
{$1000$ lần}
{$100$ lần}
{\True $10\,000$ lần}
{$5000$ lần}
\loigiai{
$50\,000 : 5 = 10\,000$ lần.\\
\textbf{Đáp án của câu hỏi là C}
}
\end{ex}

\begin{ex} %[4D3V11]
Một số có lớp nghìn là $420$ và lớp đơn vị là $035$. Số đó đọc là
\choice
{Bốn trăm hai mươi nghìn ba mươi lăm}
{\True Bốn trăm hai mươi nghìn không trăm ba mươi lăm}
{Bốn mươi hai nghìn ba mươi lăm}
{Bốn trăm hai mươi triệu ba mươi lăm}
\loigiai{
Số $420\,035$ đọc là Bốn trăm hai mươi nghìn không trăm ba mươi lăm.\\
\textbf{Đáp án của câu hỏi là B}
}
\end{ex}

\begin{ex} %[4D3V11]
Cho số $a = 687\,215$. Nếu đổi chỗ chữ số hàng trăm nghìn và chữ số hàng trăm cho nhau thì ta được số mới là
\choice
{$687\,215$}
{$867\,215$}
{$215\,687$}
{\True $287\,615$}
\loigiai{
Đổi chữ số 6 (trăm nghìn) và 2 (trăm) ta được $287\,615$.\\
\textbf{Đáp án của câu hỏi là D}
}
\end{ex}

\begin{ex} %[4D3V11]
Số gồm $5$ trăm nghìn, $4$ trăm, $3$ chục và $2$ đơn vị được viết là
\choice
{\True $500\,432$}
{$540\,032$}
{$504\,032$}
{$543\,200$}
\loigiai{
$500\,000 + 400 + 30 + 2 = 500\,432$.\\
\textbf{Đáp án của câu hỏi là A}
}
\end{ex}

\begin{ex} %[4D3V11]
Tìm số có 6 chữ số biết rằng tất cả các chữ số thuộc lớp nghìn đều bằng 3 và tổng các chữ số thuộc lớp đơn vị gấp 3 lần tổng các chữ số thuộc lớp nghìn.
\choice
{$333\,999$}
{$333\,888$}
{\True $333\,999$}
{$333\,333$}
\loigiai{
Tổng lớp nghìn: $3 + 3 + 3 = 9$.\\
Tổng lớp đơn vị: $9 \times 3 = 27 \Rightarrow 9 + 9 + 9 = 27$. Số đó là $333\,999$.\\
\textbf{Đáp án của câu hỏi là C}
}
\end{ex}

\textbf{Phần 2. Trắc nghiệm đúng sai}
\setcounterpageref{ex}{0}

\begin{ex} %[4D3V11]
Cho số $765\,432$. Xét tính đúng sai:
\choiceTF
{\True Chữ số 7 ở hàng trăm nghìn thuộc lớp nghìn.}
{\True Chữ số 4 ở hàng trăm thuộc lớp đơn vị.}
{Chữ số 6 ở hàng chục thuộc lớp đơn vị.}
{\True Lớp nghìn gồm các chữ số: $7, 6, 5$.}
\loigiai{
\begin{itemize}
    \item Chữ số 7 hàng trăm nghìn (lớp nghìn): Đúng.
    \item Chữ số 4 hàng trăm (lớp đơn vị): Đúng.
    \item Chữ số 6 ở hàng chục nghìn (lớp nghìn, không phải hàng chục lớp đơn vị). Sai.
    \item Lớp nghìn gồm: 7, 6, 5. Đúng.
\end{itemize}
\begin{itemize}
    \item a) \textbf{Đúng}.
    \item b) \textbf{Đúng}.
    \item c) \textbf{Sai}.
    \item d) \textbf{Đúng}.
\end{itemize}
}
\end{ex}

\begin{ex} %[4D3V11]
Xét cấu tạo hàng và lớp của số $804\,520$:
\choiceTF
{\True Lớp nghìn gồm $8$ trăm nghìn, $0$ chục nghìn, $4$ nghìn.}
{\True Lớp đơn vị gồm $5$ trăm, $2$ chục, $0$ đơn vị.}
{Tổng các chữ số của lớp nghìn lớn hơn tổng các chữ số của lớp đơn vị.}
{\True Tổng các chữ số của cả số là $8 + 0 + 4 + 5 + 2 + 0 = 19$.}
\loigiai{
\begin{itemize}
    \item Lớp nghìn: 804. Đúng.
    \item Lớp đơn vị: 520. Đúng.
    \item Tổng lớp nghìn $= 8 + 0 + 4 = 12 >$ Tổng lớp đơn vị $= 5 + 2 + 0 = 7$. Đúng.
    \item Tổng $= 19$. Đúng.
\end{itemize}
\begin{itemize}
    \item a) \textbf{Đúng}.
    \item b) \textbf{Đúng}.
    \item c) \textbf{Đúng}.
    \item d) \textbf{Đúng}.
\end{itemize}
}
\end{ex}

\textbf{Phần 3. Trả lời ngắn}
\setcounterpageref{ex}{0}

\begin{ex} %[4D3V11]
Trong số $482\,356$, tính tổng các chữ số thuộc lớp nghìn.
\shortans{14}
\loigiai{
\textbf{Đáp số:} $14$\\
$4 + 8 + 2 = 14$.
}
\end{ex}

\begin{ex} %[4D3V11]
Trong số $715\,489$, tính tổng các chữ số thuộc lớp đơn vị.
\shortans{21}
\loigiai{
\textbf{Đáp số:} $21$\\
$4 + 8 + 9 = 21$.
}
\end{ex}

\begin{ex} %[4D3V11]
Viết số gồm $6$ trăm nghìn, $8$ chục nghìn và $5$ đơn vị.
\shortans{680005}
\loigiai{
\textbf{Đáp số:} $680005$\\
Số đó là $680\,005$.
}
\end{ex}

\begin{ex} %[4D3V11]
Trong số $824\,562$, chữ số 2 ở hàng chục nghìn có giá trị lớn hơn chữ số 2 ở hàng đơn vị bao nhiêu đơn vị?
\shortans{19998}
\loigiai{
\textbf{Đáp số:} $19998$\\
$20\,000 - 2 = 19\,998$.
}
\end{ex}

\begin{ex} %[4D3V11]
Tìm số có 6 chữ số biết lớp nghìn là 500 và lớp đơn vị là số lớn nhất có 3 chữ số khác nhau.
\shortans{500987}
\loigiai{
\textbf{Đáp số:} $500987$\\
Số đó là $500\,987$.
}
\end{ex}

\begin{ex} %[4D3V11]
Có bao nhiêu số có 6 chữ số mà tổng các chữ số thuộc lớp nghìn bằng 2 và tổng các chữ số thuộc lớp đơn vị bằng 1?
\shortans{6}
\loigiai{
\textbf{Đáp số:} $6$\\
Lớp nghìn có tổng bằng 2: 200, 110, 101 (3 cách).\\
Lớp đơn vị có tổng bằng 1: 100, 010, 001 (3 cách).\\
Số lượng số: $3 \times 3 = 9$ số (nếu có các hoán vị 110, 101, 200).
}
\end{ex}
